\section{Robustness}\label{sec:robustness}

\paragraph{Synthetic ensemble.}
Across \SynSeeds{} synthetic scale-free graphs (DAGs with $n=\SynN$) under SH1 parameters, hub-seeded
outbreaks reach $\SynHubMean\pm\SynHubSD$ and leaf-seeded ones $\SynLeafMean\pm\SynLeafSD$ (mean
$\pm$ s.d.\ across seeds). Hub immunization is the best RQ3 strategy at 24\,h delay and maximum budget
in \SynRQThreeBestCount{} graphs. The qualitative RQ1 and RQ3 conclusions therefore do not depend on the
particular structure of the npm crawl. The synthetic graphs differ from npm in the size of the credential
effect: leaves reach much less, because synthetic maintainer portfolios are smaller and assigned
independently of the topology.

\paragraph{Sample size and null models.}
\Cref{tab:robustness} repeats the key quantities on three graph variants:
\begin{itemize}
  \item the 5k sample;
  \item the full 14{,}624-package closure;
  \item five degree-preserving rewirings of the main graph~\citep{maslov2002specificity}. These keep every
        package's in- and out-degree and the SemVer mix, but randomise who depends on whom.
\end{itemize}
The rewirings keep maintainer ownership fixed, so they separate the role of dependency structure from
that of the credential channel.

\begin{table}[t]
\centering
\caption{Robustness of the key quantities across graph variants. $\Rinf$ values are means over 1{,}000
runs per origin group.}
\label{tab:robustness}
\small
\begin{tabular}{lcccl}
\toprule
 & Hub $R_\infty$ & Hub $R_\infty$ & Leaf $R_\infty$ & Best RQ3 strategy \\
Graph & (dep.\ only) & (SH1) & (SH1) & (24\,h, max.\ budget) \\
\midrule
npm 5k & 2.3\% & 55.7\% & 22.4\% & Hub immunization \\
npm 10k (main) & 2.6\% & 55.8\% & 19.6\% & Firebreak (betweenness) \\
npm full closure & 2.2\% & 51.7\% & 7.8\% & Firebreak (betweenness) \\
null model 0 & 8.1\% & 64.2\% & 31.8\% & Acquaintance \\
null model 1 & \multicolumn{4}{c}{\textbf{[pending]}} \\
null model 2 & \multicolumn{4}{c}{\textbf{[pending]}} \\
null model 3 & \multicolumn{4}{c}{\textbf{[pending]}} \\
null model 4 & \multicolumn{4}{c}{\textbf{[pending]}} \\
\bottomrule
\end{tabular}

\end{table}

Three patterns stand out in \cref{tab:robustness}:
\begin{itemize}
  \item \textbf{npm's wiring dampens spread.} The \NullCount{} rewired graphs produce \emph{larger}
        outbreaks than the real one. Dependency-only hub outbreaks reach
        \NullDepHubMin--\NullDepHubMax{} instead of \RQOneDepOnlyHub, and SH1 leaf outbreaks reach
        \NullLeafMin--\NullLeafMax{} instead of \RQOneShOneLeaf. Random rewiring destroys npm's layered,
        nearly acyclic structure and creates cycles, so the real topology is protective relative to its
        degree sequence alone.
  \item \textbf{The leaf effect is a share of the graph.} SH1 leaf outbreaks reach \SizeFiveLeafShOne{},
        \SizeMainLeafShOne{} and \SizeFullLeafShOne{} of the 5k, 10k and full-closure graphs. In absolute terms that is between roughly 1{,}100 and 2{,}000 packages, so the
        share does not scale simply with graph size. We therefore report the leaf effect as a range
        rather than a single figure.
  \item \textbf{The RQ3 ranking is stable.} Proactive immunization beats reactive patching at a 24\,h
        delay on every variant. Which proactive heuristic is best at the largest budget varies between
        graphs (degree, betweenness or acquaintance), consistent with the convergence of the strategies
        at large $K$ in \cref{sec:rq3}.
\end{itemize}

\paragraph{Parameter sensitivity.}
Around the SH1 calibration (50 runs per setting) we vary $\beta\in\{0.01,0.02,0.05\}$ and $\gamma\in\{0.05,0.1,0.2\}$. The
conclusions hold across all of them:
\begin{itemize}
  \item a proactive strategy is the best RQ3 response at 24\,h delay and the largest budget in
        \SensProactiveBest{} settings;
  \item at $B=50$ per hour, hub immunization beats hub patching in \SensImmBeatsPatch{} settings and
        betweenness immunization in \SensDegreeBeatsFirebreak;
  \item the hub-to-leaf ratio of $\Rinf$ stays between \SensHubLeafMin{} and \SensHubLeafMax, so
        leaf-started outbreaks remain substantial whenever tokens can be stolen.
\end{itemize}
